
\documentclass{article}
\usepackage{colortbl}
\usepackage{makecell}
\usepackage{multirow}
\usepackage{supertabular}

\begin{document}

\newcounter{utterance}

\centering \large Interaction Transcript for game `wordle', experiment `medium\_frequency\_words\_no\_clue\_no\_critic', episode 14 with Qwen3.5{-}9B.
\vspace{24pt}

{ \footnotesize  \setcounter{utterance}{1}
\setlength{\tabcolsep}{0pt}
\begin{supertabular}{c@{$\;$}|p{.15\linewidth}@{}p{.15\linewidth}p{.15\linewidth}p{.15\linewidth}p{.15\linewidth}p{.15\linewidth}}
    \# & \multicolumn{2}{c}{Player} && \multicolumn{2}{c}{Game Master} \\
    \hline

    \theutterance \stepcounter{utterance}  
    & & & \multicolumn{4}{p{0.6\linewidth}}{
        \cellcolor[rgb]{0.9,0.9,0.9}{
            \makecell[{{p{\linewidth}}}]{
                \texttt{\tiny{[P1$\langle$GM]}}
                \texttt{You are a language wizard who likes to guess words by using the given rules.} \\
\\ 
\texttt{Welcome to Wordle! You have six attempts to guess the target word, a valid English word of five lowercase letters (a{-}z). Please use the tags "explanation:" and "guess:" to provide a concise explanation for each guess.} \\
\\ 
\texttt{For instance, if your guess is "apple", your response should be} \\
\texttt{explanation: this is a common five{-}letter English word, and I am starting my guess with this word.} \\
\texttt{guess: apple} \\
\\ 
\texttt{After each guess, your answer will be validated, and you will receive feedback indicating which letters are correct (green), which letters are correct but in the wrong position (yellow), and which letters are incorrect (red). This feedback can be useful in determining which letters to include or exclude in your next guess.} \\
\\ 
\texttt{For example, the feedback for "apple" might be:} \\
\texttt{guess\_feedback: a<yellow> p<yellow> p<green> l<yellow> e<red>} \\
\\ 
\texttt{The explanation should contain details about how the guess\_feedback is used to arrive at a new guess.} \\
\\ 
\texttt{Do not generate any other text.} \\
\\ 
\texttt{Let's begin with your first guess.} \\
            }
        }
    }
     \\ \\

    \theutterance \stepcounter{utterance}  
    & \multicolumn{4}{p{0.6\linewidth}}{
        \cellcolor[rgb]{0.561,0.737,0.561}{
            \makecell[{{p{\linewidth}}}]{
                \texttt{\tiny{[P1$\rangle$GM]}}
                \texttt{explanation: I am starting with "crane" because it contains common vowels (a, e) and frequent consonants (r, n) that appear in many five{-}letter words, allowing me to quickly eliminate or confirm several letters.} \\
\texttt{guess: crane} \\
            }
        }
    }
    & & \\ \\

    \theutterance \stepcounter{utterance}  
    & & & \multicolumn{2}{p{0.3\linewidth}}{
        \cellcolor[rgb]{0.9,0.9,0.9}{
            \makecell[{{p{\linewidth}}}]{
                \texttt{\tiny{[GM$|$GM]}}
                \texttt{attempts = 1} \\
\texttt{target\_word = grind} \\
\texttt{guess = crane} \\
\texttt{guess\_feedback = c<red> r<green> a<red> n<green> e<red>} \\
            }
        }
    }
    & & \\ \\

    \theutterance \stepcounter{utterance}  
    & & & \multicolumn{4}{p{0.6\linewidth}}{
        \cellcolor[rgb]{0.9,0.9,0.9}{
            \makecell[{{p{\linewidth}}}]{
                \texttt{\tiny{[P1$\langle$GM]}}
                \texttt{guess\_feedback: c<red> r<green> a<red> n<green> e<red>} \\
\\ 
\texttt{Provide your response only in this format.} \\
\texttt{explanation: details} \\
\texttt{guess: word} \\
            }
        }
    }
     \\ \\

    \theutterance \stepcounter{utterance}  
    & \multicolumn{4}{p{0.6\linewidth}}{
        \cellcolor[rgb]{0.561,0.737,0.561}{
            \makecell[{{p{\linewidth}}}]{
                \texttt{\tiny{[P1$\rangle$GM]}}
                \texttt{explanation: The feedback shows that 'r' and 'n' are correct and in the right positions (2nd and 4th), while 'c', 'a', and 'e' are incorrect. I will keep 'r' and 'n' fixed and try to find the correct 1st, 3rd, and 5th letters using common vowels and consonants.} \\
\texttt{guess: train} \\
            }
        }
    }
    & & \\ \\

    \theutterance \stepcounter{utterance}  
    & & & \multicolumn{2}{p{0.3\linewidth}}{
        \cellcolor[rgb]{0.9,0.9,0.9}{
            \makecell[{{p{\linewidth}}}]{
                \texttt{\tiny{[GM$|$GM]}}
                \texttt{attempts = 2} \\
\texttt{target\_word = grind} \\
\texttt{guess = train} \\
\texttt{guess\_feedback = t<red> r<green> a<red> i<yellow> n<yellow>} \\
            }
        }
    }
    & & \\ \\

    \theutterance \stepcounter{utterance}  
    & & & \multicolumn{4}{p{0.6\linewidth}}{
        \cellcolor[rgb]{0.9,0.9,0.9}{
            \makecell[{{p{\linewidth}}}]{
                \texttt{\tiny{[P1$\langle$GM]}}
                \texttt{guess\_feedback: t<red> r<green> a<red> i<yellow> n<yellow>} \\
\\ 
\texttt{Provide your response only in this format.} \\
\texttt{explanation: details} \\
\texttt{guess: word} \\
            }
        }
    }
     \\ \\

    \theutterance \stepcounter{utterance}  
    & \multicolumn{4}{p{0.6\linewidth}}{
        \cellcolor[rgb]{0.561,0.737,0.561}{
            \makecell[{{p{\linewidth}}}]{
                \texttt{\tiny{[P1$\rangle$GM]}}
                \texttt{explanation: The feedback confirms 'r' is correct in the 2nd position and 'n' is correct but in the wrong position (so 'n' must be in the 5th spot). 't', 'a', and 'i' are incorrect or misplaced. Since 'i' was yellow in "train" but 'n' is now known to be 5th, and 'a' is red, I need a word with 'r' at 2, 'n' at 5, and a new vowel/consonant combination. "brine" fits with 'r' at 2, 'n' at 5, and tests 'b' and 'i' (though 'i' was yellow before, maybe it's the vowel, but 'a' is red so 'i' might be the vowel). Wait, in "train", 'i' was yellow meaning it's in the word but not at 3. In "brine", 'i' is at 3. If 'i' is the vowel, let's try "brine". Actually, let's reconsider: 'r' is 2, 'n' is 5. 'a' is red. 't' is red. 'i' was yellow in "train" (t{-}r{-}a{-}i{-}n), so 'i' is in the word but not at 3. So 'i' could be at 1, 4, or 5? But 'n' is at 5. So 'i' is at 1 or 4. Let's try "brine" {-}> b{-}r{-}i{-}n{-}e. 'e' is new. Or "crine"? No. How about "prone"? p{-}r{-}o{-}n{-}e. 'o' is new. Let's go with "prone" to test 'p' and 'o'.} \\
\texttt{guess: prone} \\
            }
        }
    }
    & & \\ \\

    \theutterance \stepcounter{utterance}  
    & & & \multicolumn{2}{p{0.3\linewidth}}{
        \cellcolor[rgb]{0.9,0.9,0.9}{
            \makecell[{{p{\linewidth}}}]{
                \texttt{\tiny{[GM$|$GM]}}
                \texttt{attempts = 3} \\
\texttt{target\_word = grind} \\
\texttt{guess = prone} \\
\texttt{guess\_feedback = p<red> r<green> o<red> n<green> e<red>} \\
            }
        }
    }
    & & \\ \\

    \theutterance \stepcounter{utterance}  
    & & & \multicolumn{4}{p{0.6\linewidth}}{
        \cellcolor[rgb]{0.9,0.9,0.9}{
            \makecell[{{p{\linewidth}}}]{
                \texttt{\tiny{[P1$\langle$GM]}}
                \texttt{guess\_feedback: p<red> r<green> o<red> n<green> e<red>} \\
\\ 
\texttt{Provide your response only in this format.} \\
\texttt{explanation: details} \\
\texttt{guess: word} \\
            }
        }
    }
     \\ \\

    \theutterance \stepcounter{utterance}  
    & \multicolumn{4}{p{0.6\linewidth}}{
        \cellcolor[rgb]{0.561,0.737,0.561}{
            \makecell[{{p{\linewidth}}}]{
                \texttt{\tiny{[P1$\rangle$GM]}}
                \texttt{explanation: The feedback confirms 'r' is correct in the 2nd position and 'n' is correct in the 4th position. 'p', 'o', and 'e' are incorrect. Since 'i' was yellow in "train" (meaning it is in the word but not at position 3), and 'n' is now confirmed at position 4, 'i' must be at position 1 or 5. However, 'e' is red, so 'i' cannot be at 5 if the word ends in 'e'. Wait, in "train", 'n' was yellow, meaning it was in the word but not at 5. Now in "prone", 'n' is green at 4. So 'n' is definitely at 4. 'i' was yellow in "train" at position 3, so 'i' is in the word but not at 3. With 'n' at 4, 'i' could be at 1 or 5. But 'e' is red, so the word doesn't end in 'e'. Let's try "irony": i{-}r{-}o{-}n{-}y. 'o' is red, so that's out. How about "irons"? i{-}r{-}o{-}n{-}s. 'o' is red. Maybe "inlet"? i{-}n{-}l{-}e{-}t. 'n' would be at 2, but we need 'n' at 4. Let's reconsider the letters: \_ r \_ n \_. 'i' is in the word, not at 3. 'a', 'e', 'o', 't', 'c', 'p' are out. 'i' could be at 1 or 5. If 'i' is at 1: i r \_ n \_. If 'i' is at 5: \_ r \_ n i. Let's try "irons" again? No, 'o' is red. What about "inrun"? Not a word. "inert"? 'e' is red. "inlet"? 'e' is red. Maybe "inlay"? i{-}n{-}l{-}a{-}y. 'n' at 2, no. Let's try "inlet" logic again. Wait, 'i' was yellow in "train" (t{-}r{-}a{-}i{-}n). 'i' is at 3 in "train", so it's not at 3. In the target, 'n' is at 4. So 'i' is at 1 or 5. Let's try "inlet" {-}> no, 'n' at 2. Try "inrun" {-}> no. How about "inlet" is wrong. Let's try "inlet" is not it. What words have \_ r \_ n \_? "arson", "baron", "carbon", "drown", "frown",} \\
            }
        }
    }
    & & \\ \\

    \theutterance \stepcounter{utterance}  
    & & & \multicolumn{2}{p{0.3\linewidth}}{
        \cellcolor[rgb]{0.9,0.9,0.9}{
            \makecell[{{p{\linewidth}}}]{
                \texttt{\tiny{[GM$|$GM]}}
                \texttt{Error: The response should contain the 'guess:' keyword exactly once.} \\
            }
        }
    }
    & & \\ \\

    \theutterance \stepcounter{utterance}  
    & & & \multicolumn{2}{p{0.3\linewidth}}{
        \cellcolor[rgb]{0.9,0.9,0.9}{
            \makecell[{{p{\linewidth}}}]{
                \texttt{\tiny{[GM$|$GM]}}
                \texttt{game\_result = ABORT} \\
            }
        }
    }
    & & \\ \\

\end{supertabular}
}

\end{document}
